\chapter{Discussion}

\section{Performance (Research Question 1)}

The first research question asks how the two algorithms compare in execution
time and throughput across message sizes on each platform.
The measurements give a clear answer: Ascon-AEAD128 was faster than AES-128-GCM
at every measured message size, on both boards, and at every optimization level.
The advantage was present in both the fixed per-call cost that dominates short
messages (Section~\ref{sec:exec-time}) and the per-byte cost that dominates long
messages (Section~\ref{sec:throughput}).
Ascon-AEAD128 leads in both regimes, so the conclusion does not depend on the
message-size profile of the application.

The reason for the difference lies in the structure of the two algorithms.
AES in software relies on table lookups for the substitution step, and the GCM
authentication step performs a multiplication in a binary field, built from many
smaller operations when no wide hardware multiplier is available.
Ascon-AEAD128 is built from a permutation of bitwise operations, exclusive-or,
bitwise-and, and rotation, applied to a small state held in registers.
These operations are native to both cores and require no table memory.
The single-pass duplex construction also avoids the separate encryption and
authentication passes that AES-128-GCM performs.
The measured advantage is the practical consequence of this design difference.

\section{Flash Footprint (Research Question 2)}

The second research question asks about the Flash footprint of each
implementation relative to the constraints of a small microcontroller.
The answer depends on the optimization level, and the size-optimized level is
the one relevant to a memory-constrained deployment.
At \texttt{-Os}, AES-128-GCM required approximately 3.8 times the code space of
Ascon-AEAD128, consistently on both boards (Section~\ref{sec:code-size}).
The small code size of Ascon-AEAD128 follows from its single permutation and its
lack of precomputed lookup tables, whereas AES-128-GCM brings the AES tables
and the field-multiplication code.

This advantage at \texttt{-Os} does not hold at the speed-oriented optimization
levels, for the reason examined in Section~\ref{sec:optlevel}.
The footprint result is therefore reported at the optimization level that a
constrained deployment would actually choose, where Ascon-AEAD128 is both
smaller and, from the timing results, faster.

\section{Optimization Sensitivity (Research Question 3)}
\label{sec:optlevel}

The third research question asks how the compiler optimization level affects the
relative performance of the two algorithms.
The measurements show that the optimization level matters a great deal, in two
distinct ways.

The first concerns execution time.
The largest single improvement for both algorithms came from enabling
optimization at all --- the step from \texttt{-O0} to \texttt{-O1}.
At 1000 bytes on the Cortex-M0 this step reduced the Ascon-AEAD128 time by a
factor of about eight and the AES-128-GCM time by a factor of about three.
This uneven response means the relative advantage of Ascon-AEAD128 is
understated by any measurement taken at \texttt{-O0}.
A study that benchmarked only unoptimized code would have reported the two
algorithms as close --- with a ratio of 1.16 on the Cortex-M4F --- and would
have missed the larger advantage that appears under the optimization a real
build would use.
The choice to sweep the optimization level, rather than report a single level,
is justified by this observation.

The behavior at the higher levels was not uniform.
On the Cortex-M4F, AES-128-GCM gained nothing from \texttt{-O3} over
\texttt{-O2}, consistent with the common observation that aggressive
optimization yields little for table-driven code on these cores, while
Ascon-AEAD128 continued to improve at \texttt{-O3}.
On the Cortex-M0 the pattern was reversed, with AES-128-GCM improving
substantially at \texttt{-O3} and Ascon-AEAD128 improving only modestly.
The \texttt{-Os} level was slower than \texttt{-O2} for both algorithms on both
boards, because size optimization disables the inlining and loop unrolling that
the cryptographic inner loops benefit from.

The second way the optimization level matters concerns code size, and it is the
more striking effect.
The Flash cost of the Ascon-AEAD128 reference implementation varied strongly
with the optimization level, while the AES-128-GCM cost was comparatively
stable.
The Ascon-AEAD128 cost rose from \texttt{-O0} to a peak at \texttt{-O3},
reaching 43{,}024 bytes on the Cortex-M0 and 34{,}560 bytes on the Cortex-M4F,
then fell sharply at \texttt{-Os} to roughly 4\,KB.
The cause is the unrolling of the permutation: at the speed-oriented levels the
compiler unrolls the permutation rounds into straight-line code, which the
authenticated-encryption routine includes at several call sites, expanding the
binary.
At \texttt{-Os} the compiler keeps the rounds rolled, and the code collapses to
a small size.
The same shape appeared on both boards, showing that the effect is a property
of the implementation and the compiler rather than of the processor.

This links the second and third research questions.
The footprint advantage of Ascon-AEAD128 is real but conditional on the
optimization level: it exists at \texttt{-Os} and is reversed at the
speed-oriented levels.
The advantage and this condition should be reported together.

\section{Architecture Dependence (Research Question 4)}

The fourth research question asks whether the performance gap between the two
algorithms differs between the Cortex-M0 and the Cortex-M4F, and what this
implies for device selection.
The measurements show that the gap is consistently larger on the weaker core.

At equal clock frequency and at 1000 bytes, the ratio of AES-128-GCM time to
Ascon-AEAD128 time was about 4.6 at \texttt{-O1} and \texttt{-O2} on the
Cortex-M0, against about 2.0 at the same levels on the Cortex-M4F.
At every optimized level the ratio was larger on the Cortex-M0.
The advantage of Ascon-AEAD128 therefore grows as the processor becomes more
constrained.

This can be read against the architectural baseline established in the harness
validation.
For the plain memory workload the Cortex-M0 was about 2.8 times slower than
the Cortex-M4F at the same clock.
If both algorithms simply scaled with this architectural factor, their ratio to
each other would be the same on both boards.
The fact that the AES-128-GCM disadvantage grows by more than this factor on
the Cortex-M0 indicates that AES-128-GCM contains operations that are
especially costly on the weaker core.
The field multiplication in the GCM authentication step is the principal
example: it must be synthesized from many small operations when no wide
hardware multiplier is present, which is the case on the Cortex-M0.
Ascon-AEAD128, built from operations the Cortex-M0 executes natively, stays
closer to the architectural baseline.

The implication for device selection is direct: the more constrained the
processor, the stronger the case for Ascon-AEAD128 over AES-128-GCM.
On the lowest-end cores, which are also the most numerous and the most
cost-sensitive, the advantage is largest.

\section{The Fairness of the Comparison}

The comparison was designed to avoid handicapping either algorithm, and one
configuration choice deserves explicit defense.
AES-128-GCM was measured with the small field-multiplication table, the
memory-frugal configuration disclosed in Section~3.4.1.
The larger table would improve throughput at a RAM cost that Ascon-AEAD128 does
not require, and on a part with 16\,KB of SRAM that is exactly the kind of cost
a constrained deployment may decline.

The reported AES-128-GCM figures therefore represent the memory-frugal end of
AES-128-GCM rather than its fastest possible configuration.
Enabling the larger table would narrow the throughput gap at large messages, at
a RAM cost Ascon-AEAD128 does not pay, and would not change the code-size or
small-message conclusions.
The comparison is thus a fair one between two lightweight configurations, and
the configuration choice is named and justified so the result cannot be read as
a weakened AES-128-GCM.

\section{Practical Conclusions}

The results support a consistent practical recommendation for the class of
devices studied.
For authenticated encryption on a constrained microcontroller without hardware
AES acceleration, Ascon-AEAD128 is the stronger choice on the measured criteria:
it was faster than AES-128-GCM at every message size and optimization level,
smaller in code at the size-optimized level that such a device would use, and
its advantage was largest on the most constrained processor.

Two qualifications accompany this recommendation.
First, the speed advantage and the size advantage do not occur at the same
optimization level for the reference implementation.
The size advantage holds at \texttt{-Os}, where Ascon-AEAD128 is both smaller
and faster than AES-128-GCM, so a deployment that builds for size obtains both
benefits at once.
A deployment that builds for speed at \texttt{-O2} or \texttt{-O3} obtains the
speed advantage but accepts a larger Ascon-AEAD128 code size due to the
unrolling of the permutation.
Second, these conclusions concern software-only implementations.
On a device with a hardware AES accelerator the balance would differ, and that
case was deliberately outside the scope of this work in order to isolate the
software comparison.

Within those qualifications, the measurements give a quantitative basis for
preferring Ascon-AEAD128 on constrained, software-only targets, and they
quantify how the advantage grows as the device becomes more limited.
